\section*{Chapitre \ref{ch:g8:ice-on-earth} --- La glace sur la Terre}

\begin{solution}{exo:g8:ice-on-earth:1}
(a) La zone d'accumulation. (b) La zone d'ablation. (c) Le front
glaciaire. (d) Une \omterm{def:g8:ice-on-earth:glacial-landforms}{moraine} (frontale).
\end{solution}

\begin{solution}{exo:g8:ice-on-earth:2}
(a) Faux~: un \omterm{def:g8:ice-on-earth:glacier}{glacier} est fait de \omterm{def:g4:clouds-rain-snow:precipitation}{neige} tombée sur les terres et
transformée en glace~; l'eau de \omterm{def:g2:water-on-earth:ocean}{mer} gelée, c'est la \omterm{def:g8:ice-on-earth:sea-ice}{banquise}. (b) Vrai.
(c) Faux~: sa glace continue de s'écouler vers le bas~; le front recule
parce que la glace fond plus vite qu'elle n'avance. (d) Vrai.
\end{solution}

\begin{solution}{exo:g8:ice-on-earth:3}
(a) Une \omterm{def:g8:ice-on-earth:glacial-landforms}{vallée en U}. (b) Un bloc erratique. (c) Une \omterm{def:g8:ice-on-earth:ice-sheet}{plate-forme de glace}. (d) Le \omterm{def:g8:ice-on-earth:permafrost}{pergélisol}.
\end{solution}

\begin{solution}{exo:g8:ice-on-earth:4}
(a) L'Antarctique et le Groenland. (b) L'Antarctique, d'environ
\qty{58}{m} (le Groenland~: environ \qty{7.4}{m}). (c) Environ un dixième
(10~\%).
\end{solution}

\begin{solution}{exo:g8:ice-on-earth:5}
(a) Environ \num{7.05}~millions de~\unit{km^2} en 1979 et
\num{3.57}~millions de~\unit{km^2} en 2012. (b) $7.05 - 3.57 = 3.48$~:
environ \num{3.5}~millions de~\unit{km^2}. (c) $3.48 \div 7.05 \approx 0.49$~:
environ la moitié. (d) Une année isolée dépend du temps qu'il a fait cet
été-là~: 2012 est bien au-dessous de la droite. La droite de tendance
utilise toutes les années, si bien qu'un été inhabituel la modifie très
peu.
\end{solution}

\begin{solution}{exo:g8:ice-on-earth:6}
(a) $2004 - 1941 = 63$ ans. (b) $12000 \div 63 \approx 190.5$~: environ
\qty{190}{m} par an. (c) $2021 - 1984 = 37$ ans~; $3200 \div 37 \approx
86$~: environ \qty{86}{m} par an. (d) Le \omterm{def:g8:ice-on-earth:glacier}{glacier} Muir, plus de deux fois
plus vite.
\end{solution}

\begin{solution}{exo:g8:ice-on-earth:7}
Dans la \emph{zone d'accumulation}, haut sur la \omterm{def:g1:landscapes-around-us:landform}{montagne}, il tombe chaque
année plus de \omterm{def:g4:clouds-rain-snow:precipitation}{neige} qu'il n'en fond. La \omterm{def:g4:clouds-rain-snow:precipitation}{neige} nouvelle recouvre et tasse
l'ancienne, qui se change en \emph{névé} au bout d'environ un an, puis,
plus profondément et bien plus tard, en glace compacte. Sous son propre
poids, la glace \emph{s'écoule} lentement vers le bas de la \omterm{def:g1:landscapes-around-us:landform}{vallée}. Plus
bas, dans la \emph{zone d'ablation}, elle fond plus vite que la \omterm{def:g4:clouds-rain-snow:precipitation}{neige} ne
tombe~; la glace qui atteint le front s'est formée tout en haut, souvent
il y a longtemps.
\end{solution}

\begin{solution}{exo:g8:ice-on-earth:8}
(a) Dans le verre B~: la glace fondue coule dans l'eau et en fait monter
le niveau~; dans le verre A, le niveau reste le même. (b) Le verre A est
l'\omterm{def:g8:ice-on-earth:ice-sheet}{iceberg} (de la glace qui flotte déjà dans la \omterm{def:g2:water-on-earth:ocean}{mer})~; le verre B est la
\omterm{def:g8:ice-on-earth:ice-sheet}{calotte glaciaire} (de la glace sur les terres). (c) Une \omterm{def:g8:ice-on-earth:ice-sheet}{plate-forme de
glace} retient les \omterm{def:g8:ice-on-earth:glacier}{glaciers} situés derrière elle. Quand elle se disloque,
ces \omterm{def:g8:ice-on-earth:glacier}{glaciers} accélèrent et déversent davantage de glace terrestre dans la
\omterm{def:g2:water-on-earth:ocean}{mer}, ce qui la fait monter.
\end{solution}

\begin{solution}{exo:g8:ice-on-earth:9}
(a) Une \omterm{def:g8:ice-on-earth:glacial-landforms}{vallée en U}, creusée par un \omterm{def:g8:ice-on-earth:glacier}{glacier}. (b) Un \omterm{def:g8:ice-on-earth:glacier}{glacier} l'a apporté
de loin en amont dans la \omterm{def:g1:landscapes-around-us:landform}{vallée}, là où affleure du granite, et l'a
abandonné quand la glace a fondu~: c'est un bloc erratique. (c) Une
\omterm{def:g8:ice-on-earth:glacial-landforms}{moraine} latérale~: sa hauteur montre l'épaisseur qu'avait autrefois le
\omterm{def:g8:ice-on-earth:glacier}{glacier}~; la glace montait au moins jusque-là sur le versant.
\end{solution}

\begin{solution}{exo:g8:ice-on-earth:10}
(a) $7.4 \div 10 = \qty{0.74}{m}$. (b) $58 \div 100 = \qty{0.58}{m}$.
(c) $58 \div 7.4 \approx 7.8$~: environ 8 fois plus.
\end{solution}

\begin{solution}{exo:g8:ice-on-earth:11}
(a) $7.21 \times 10^6 = \qty{7210000}{km^2}$ et $4.65 \times 10^6 =
\qty{4650000}{km^2}$. (b) $7.21 - 4.65 = \num{2.56}\text{ millions de }\unit{km^2}$, soit
\qty{2560000}{km^2}. (c) $2.56 \div 3 \approx 0.85$~: environ
\num{0.85}~million de~\unit{km^2} par décennie, un peu plus que les
\num{0.76}~million de~\unit{km^2} de la droite de tendance~; la moyenne de
chaque décennie porte encore les hauts et les bas de ses dix étés, mais
les deux façons de lire le graphique montrent la même forte perte. (d)
L'étendue ne peut pas descendre au-dessous de zéro, et une droite de
tendance ne décrit que les années qui ont été mesurées~: la suite dépend
de l'évolution du \omterm{def:g5:weather-and-climate:climate}{climat} de l'Arctique, si bien que la droite n'est pas
une prédiction.
\end{solution}

\begin{solution}{exo:g8:ice-on-earth:12}
(a) $120 \div 1 = 120$~: la \omterm{def:g1:landscapes-around-us:water}{côte} se trouvait environ \qty{120}{km} plus
loin vers le large. (b) Elle s'était évaporée de l'\omterm{def:g2:water-on-earth:ocean}{océan}, était tombée en
\omterm{def:g4:clouds-rain-snow:precipitation}{neige} sur les \omterm{def:g3:continents-and-oceans:continent}{continents} et y était restée sous forme de glace, dans les
grandes \omterm{def:g8:ice-on-earth:ice-sheet}{calottes glaciaires}. (c) Une \omterm{def:g8:ice-on-earth:glacial-landforms}{moraine} frontale se forme à
l'extrémité d'un \omterm{def:g8:ice-on-earth:glacier}{glacier}. Une \omterm{def:g8:ice-on-earth:glacial-landforms}{moraine} dans une \omterm{def:g1:landscapes-around-us:landform}{plaine} loin de toute
\omterm{def:g1:landscapes-around-us:landform}{montagne} montre qu'une \omterm{def:g8:ice-on-earth:ice-sheet}{calotte glaciaire} a autrefois atteint cette
\omterm{def:g1:landscapes-around-us:landform}{plaine}, à une époque où la glace couvrait bien plus de terres
qu'aujourd'hui.
\end{solution}

\begin{solution}{pb:g8:ice-on-earth:1}
\textbf{1.} $1080 - 0 = \qty{1080}{m}$.

\textbf{2.} 1960--1970~: \qty{180}{m}~; 1970--1980~: \qty{150}{m}~;
1980--1990~: \qty{190}{m}~; 1990--2000~: \qty{180}{m}~; 2000--2010~:
\qty{180}{m}~; 2010--2020~: \qty{200}{m}. Le moins en 1970--1980, le plus
en 2010--2020.

\textbf{3.} Non~: un \omterm{def:g8:ice-on-earth:glacier}{glacier} qui recule continue de s'écouler vers le bas.
Son front recule parce que la glace fond plus vite qu'elle n'arrive.

\textbf{4.} Une \omterm{def:g8:ice-on-earth:glacial-landforms}{moraine} frontale~: la glace a transporté des débris
rocheux vers le bas de la \omterm{def:g1:landscapes-around-us:landform}{vallée} et les a déposés à son front, là où la
glace fondait, pendant que le front se tenait à cet endroit.

\textbf{5.} $1080 \div 6 = \qty{180}{m}$ par décennie.

\textbf{6.} $180 \div 10 = \qty{18}{m}$ par an.

\textbf{7.} $190 \div 18 \approx 10.6$~: environ 11 fois plus vite.

\textbf{8.} Moins~: il a perdu plus de glace dans sa zone d'ablation qu'il
n'en a gagné dans sa zone d'accumulation, si bien que son front a reculé.

\textbf{9.} En U~: un fond large aux versants raides, parce que le \omterm{def:g8:ice-on-earth:glacier}{glacier}
a élargi et approfondi la \omterm{def:g1:landscapes-around-us:landform}{vallée} pendant qu'il la remplissait.

\textbf{10.} $1980 - 1080 = \qty{900}{m}$.

\textbf{11.} $900 \div 18 = 50$ ans~: en \textbf{2070}.

\textbf{12.} La vitesse a changé d'une décennie à l'autre (de 150 à
\qty{200}{m}), et le recul futur dépend des chutes de \omterm{def:g4:clouds-rain-snow:precipitation}{neige} et des
températures estivales des années à venir~: une vitesse passée n'est pas
une prédiction.
\end{solution}
